% Akronimy używane w pracy magisterskiej
\newacronym{DC}{DC}{Direct Current (prąd stały)}
\newacronym{AC}{AC}{Alternating Current (prąd przemienny)}
\newacronym{OZE}{OZE}{Odnawialne Źródła Energii}
\newacronym{NMP}{NMP}{Non-Minimum Phase (faza nieminimalna)}
\newacronym{RHP}{RHPZ}{Right Half-Plane Zero (zero w prawej półpłaszczyźnie zmiennej s)}
\newacronym{MFBB}{MF-BB}{Model-Free Bang-Bang (bezmodelowe sterowanie przekaźnikowe)}
\newacronym{TrendMFBB}{Trend MF-BB}{Model-Free Bang-Bang z filtrowanym trendem napięcia}
\newacronym{FSMPC}{FS-MPC}{Finite-State Model Predictive Control (sterowanie predykcyjne o skończonym zbiorze stanów)}
\newacronym{PSO}{PSO}{Particle Swarm Optimization (Optymalizacja Rojem Cząstek)}
\newacronym{ZOH}{ZOH}{Zero-Order Hold (ekstrapolator zerowego rzędu)}
\newacronym{OCP}{OCP}{Over-Current Protection (zabezpieczenie nadprądowe)}
\newacronym{PWM}{PWM}{Pulse Width Modulation (modulacja szerokości impulsów)}
\newacronym{PI}{PI}{Proportional-Integral (regulator proporcjonalno-całkujący)}
\newacronym{LPF}{LPF}{Low-Pass Filter (filtr dolnoprzepustowy)}
\newacronym{ESR}{ESR}{Equivalent Series Resistance (zastępcza rezystancja szeregowa)}
\newacronym{DT}{DT}{Digital Twin (Bliźniak Cyfrowy)}
\newacronym{OOP}{OOP}{Object-Oriented Programming (programowanie obiektowe)}
\newacronym{IAE}{IAE}{Integral of Absolute Error (całka modułu błędu)}
\newacronym{ISE}{ISE}{Integral of Squared Error (całka kwadratu błędu)}
\newacronym{ITAE}{ITAE}{Integral of Time-weighted Absolute Error (całka iloczynu czasu i modułu błędu)}
\newacronym{MSE}{MSE}{Mean Squared Error (błąd średniokwadratowy)}
\newacronym{PW}{PW}{Politechnika Warszawska}
\newacronym{IEEE}{IEEE}{Institute of Electrical and Electronics Engineers}

% Symbole i oznaczenia matematyczne
\newglossaryentry{symb:Vin}{
    name={\ensuremath{V_{\mathrm{in}}}},
    description={Napięcie zasilania przekształtnika (strona niskonapięciowa) [V]},
    sort=Vin, type=symbols
}
\newglossaryentry{symb:L}{
    name={\ensuremath{L}},
    description={Indukcyjność dławika magazynującego energię [H]},
    sort=L, type=symbols
}
\newglossaryentry{symb:RL}{
    name={\ensuremath{R_L}},
    description={Pasożytnicza rezystancja szeregowa uzwojenia dławika [\textOmega]},
    sort=RL, type=symbols
}
\newglossaryentry{symb:C}{
    name={\ensuremath{C}},
    description={Pojemność kondensatora filtru wyjściowego [F]},
    sort=C, type=symbols
}
\newglossaryentry{symb:Rload}{
    name={\ensuremath{R_{\mathrm{load}}}},
    description={Rezystancja zastępcza obciążenia wyjściowego [\textOmega]},
    sort=Rload, type=symbols
}
\newglossaryentry{symb:iL}{
    name={\ensuremath{i_L(t)}},
    description={Chwilowy prąd dławika przekształtnika [A]},
    sort=iL, type=symbols
}
\newglossaryentry{symb:udc}{
    name={\ensuremath{u_{\mathrm{dc}}(t)}},
    description={Chwilowe napięcie wyjściowe na kondensatorze [V]},
    sort=udc, type=symbols
}
\newglossaryentry{symb:uref}{
    name={\ensuremath{u_{\mathrm{ref}}}},
    description={Wartość zadana napięcia wyjściowego przekształtnika [V]},
    sort=uref, type=symbols
}
\newglossaryentry{symb:s}{
    name={\ensuremath{s(t)}},
    description={Dwustanowy sygnał sterujący kluczami półprzewodnikowymi, $s \in \{0, 1\}$},
    sort=s, type=symbols
}
\newglossaryentry{symb:D}{
    name={\ensuremath{D}},
    description={Współczynnik wypełnienia impulsów sterujących (ang. duty cycle)},
    sort=D, type=symbols
}
\newglossaryentry{symb:ides}{
    name={\ensuremath{i_{\mathrm{des}}}},
    description={Pożądany (zadany) prąd dławika wyznaczany w regulatorze MF-BB [A]},
    sort=ides, type=symbols
}
\newglossaryentry{symb:wi}{
    name={\ensuremath{w_i}},
    description={Współczynnik wagowy uchybu prądu w funkcji przełączającej [-]},
    sort=wi, type=symbols
}
\newglossaryentry{symb:wu}{
    name={\ensuremath{w_u}},
    description={Współczynnik wagowy trendu napięcia wyjściowego [-]},
    sort=wu, type=symbols
}
\newglossaryentry{symb:delta_uf}{
    name={\ensuremath{\Delta u_f}},
    description={Przyrost przefiltrowanego napięcia wyjściowego w kroku regulatora [V]},
    sort=duf, type=symbols
}
\newglossaryentry{symb:fci}{
    name={\ensuremath{f_{c,i}}},
    description={Częstotliwość odcięcia filtru dolnoprzepustowego prądu [Hz]},
    sort=fci, type=symbols
}
\newglossaryentry{symb:fcu}{
    name={\ensuremath{f_{c,u}}},
    description={Częstotliwość odcięcia filtru dolnoprzepustowego napięcia [Hz]},
    sort=fcu, type=symbols
}
\newglossaryentry{symb:Ts}{
    name={\ensuremath{T_s}},
    description={Okres próbkowania i taktowania algorytmu sterowania ($10~\mu\mathrm{s}$) [s]},
    sort=Ts, type=symbols
}
\newglossaryentry{symb:dt}{
    name={\ensuremath{\Delta t}},
    description={Krok całkowania numerycznego silnika fizyki ($0{,}5~\mu\mathrm{s}$) [s]},
    sort=dt, type=symbols
}
\newglossaryentry{symb:w_inertia}{
    name={\ensuremath{w}},
    description={Waga bezwładności w algorytmie PSO (ang. inertia weight) [-]},
    sort=w_inertia, type=symbols
}
\newglossaryentry{symb:c1}{
    name={\ensuremath{c_1}},
    description={Współczynnik uczenia kognitywnego cząstki w roju PSO [-]},
    sort=c1, type=symbols
}
\newglossaryentry{symb:c2}{
    name={\ensuremath{c_2}},
    description={Współczynnik uczenia społecznego roju w PSO [-]},
    sort=c2, type=symbols
}
